% Lecture 10, page 1 -- manuscript dated 6-XII-1994.
% Crossed-out intermediate calculations in the manuscript are omitted.

\clearpage
\section*{Lecture 10 - Page 1}

\begin{center}
\underline{\textbf{ALGEBRA}}\\[2pt]
\textit{(LECTURE 10)}
\end{center}
\begin{flushright}
\textit{6-Dec-1994}
\end{flushright}

\bigskip
\noindent\underline{\textbf{Action of a group on a set. Sylow's Theorems 1, 2, 3.}}

\bigskip
\noindent\underline{Def.} Let $(G,\cdot)$ be a group, and let $A\neq\varnothing$ be a set.

\noindent A map $G\times A\xrightarrow{\ \psi\ } A$,
\quad $(a,x)\overset{\psi}{\longmapsto}a\cdot x\in A$, with the properties:

\bigskip
\noindent (1)\quad $a\cdot(b\cdot x)=(a\cdot b)\cdot x$,
\quad $(\forall)\,a,b\in G$, \quad $(\forall)\,x\in A$.

\medskip
\noindent (2)\quad $e\cdot x=x$.

\bigskip
\noindent\underline{Ex.} \textbf{1)} Let $G$ be an arbitrary group, $A=G$,
\quad $G\times G\xrightarrow{\ \psi\ }G$.

\noindent $(a,x)\overset{\psi}{\longmapsto}a\cdot x=axa^{-1}$
(the conjugate of $x$ by $a$).

\bigskip
\noindent (1)\quad $a\cdot(b\cdot x)=a(bxb^{-1})a^{-1}=abxb^{-1}a^{-1}$.

\medskip
\begin{center}
\fbox{$\begin{aligned}
 a\cdot(b\cdot x)&=(ab)\cdot x,\\
 e\cdot x&=x.
\end{aligned}$}
\end{center}

\noindent $\psi$ -- the action by conjugation.

\bigskip
\noindent\textbf{2)} $G$ a group, \quad $A=\{H\mid H\leq G\}$.

\noindent Define $G\times A\xrightarrow{\ \psi\ }A$,
\quad $(g,H)\overset{\psi}{\longmapsto}g\cdot H=gHg^{-1}\in A$.

\bigskip
\noindent $f_a:G\longrightarrow G$, \quad $f_a(x)=axa^{-1}$,
\quad $f_a\in\operatorname{Aut}(G)$.

\bigskip
\noindent $f_a(xy)=axya^{-1}$.

\noindent $f_a(x)\cdot f_a(y)=axa^{-1}aya^{-1}=axya^{-1}$.

\bigskip
\noindent $f_a(x_1)=f_a(x_2)\ \Longrightarrow\ ax_1a^{-1}=ax_2a^{-1}
\ \Longrightarrow\ x_1=x_2$.

\noindent $\Longrightarrow f_a$ is injective.

\medskip
\noindent $(\forall)\,y\in G$, $(\exists)\,x\in G$ such that $f_a(x)=y$;
\quad $x=a^{-1}ya$ $\Longrightarrow f_a$ is surjective.

\bigskip
\noindent $\Longrightarrow f_a\in\operatorname{Aut}(G)$.
\quad $f_a$ is called an \underline{inner automorphism} of $G$.

\bigskip
\noindent $aHa^{-1}=f_a(H)\in A$.

\bigskip
\noindent $\psi$ satisfies (1) and (2). $\psi$ is called
\underline{the conjugation action} of $G$ on the set of its subgroups.
